命题 \ref{L:trace cat} 的证明中定义的元素 $S_{u,v}$ 恰好映到构造 \ref{Con: S} 中构造的 $S$-算子。特别地，元素 \eqref{E:abs S} 映到 \eqref{E: SLa}。

\begin{ex}
现在把上述一般讨论具体到几个特殊情形。
我们考虑 $G$ 约化、$K=B$ 的情形。回忆 $\on{D}_c^b(B\backslash G/B)$ 通常称为有限 Hecke 范畴。
若 $E=\Ql$ 且 $k$ 是有限域，则对 $w\in W$，记 $\IC_w$ 为 $B\backslash BwB/B$ 上局部系统 $\Ql[\ell(w)](\ell(w)/2)$ 到其闭包的中间延拓，其中 $\ell(w)=\dim BwB/B$ 是 $W$ 上通常的长度函数。（这里我们固定了半 Tate 扭转。）
于是 $\on{D}_c^b(B\backslash G/B)$ 中由 $\{\IC_w[d](\frac{d}{2}),\ w\in W, d\in \bZ\}$ 张成的子范畴是半单幺半子范畴。它的 Grothendieck 环同构于 $\bZ[v,v^{-1}]$ 上的有限 Hecke 代数 $\mathbb H_W$（其中 $\bZ[v,v^{-1}]\subset \mathbb H_W$ 的嵌入由 $\{\on{IC}_0(i), \ i\in\bZ\}$ 到 Hecke 范畴的嵌入诱导）。当 $v=1$ 时，$\mathbb H_W$ 特化为 $W$ 的群代数；当 $v=\sqrt{q}$ 时，它特化为前面提到的 Hecke 代数 $H_W$。
由引理 \ref{L: K iso End} 可知
\begin{itemize}
\item 在例 \ref{Ex:springer} 的情形，$\delta_B$ 上有 $\bH^{\on{op}}_W$ 的自然作用。但显然 $S_{\IC_w(i)}=S_{\IC_w}$。因此 $\bH^{\on{op}}_W$ 的作用经过 $W$ 在 $\delta_B$ 上的作用分解。事实上，我们的构造不过是把 Ginzburg 用对应构造 Springer 作用的做法重新表述了一遍。
\item 在例 \ref{Ex:DL} 的情形，映射 \eqref{E: K to End3} 诱导 $H_W^{\on{op}}$ 在 $\widetilde{\delta_e}$ 上的典范作用。另一方面，
由 \eqref{E:Hk for unit} 有 $\End\widetilde{\delta_e}=H_W^{\on{op}}$。于是得到态射 $H^{\on{op}}_W\to H^{\on{op}}_W$。下面的命题 \ref{P:S=T finite} 说明该映射是恒等映射。
\end{itemize}
\end{ex}
\begin{rmk}
映射 \eqref{E: K to End2}（及其 $\sigma$-扭版本）应当使人能够推断 $\on{D}^C(X)$ 中更一般对象（例如 Deligne–Lusztig 表示）的自同态环的信息。
\end{rmk}

设 $\mF\in \on{D}_c^b(B\backslash G/B)$，并设它带有典范同构 $\sigma^*\mF\cong \mF$（Weil 结构）。记 $f_\mF$ 为经由 Grothendieck 的层-函数字典得到的 $B(\bF_q)\backslash G(\bF_q)/B(\bF_q)$ 上的函数。由 \eqref{E:Hk for unit}，该函数给出 $\End(\widetilde{\delta_e})$ 中的元素，记作 $S'_\mF$。
\begin{prop}
\label{P:S=T finite}
在映射 $\End_{\Tr_{\sigma}(\on{D}^b_c(B\backslash G/B))}(\widetilde{\mathbf{1}})\to \End (\widetilde{\delta}_e)$ 下，元素 $S_\mF$ 映到 $S'_{\mF}$。
\end{prop}
\begin{proof}
我们给出该命题的详细证明。
\cite[\S 6.3]{XZ} 中用了一个类似但更复杂的论证来建立它的仿射版本（该版本将在下文作为定理 \ref{T: trace of geom Sat} 的一部分陈述）。

记 $\mZ=m^{-1}(\frac{B}{\Ad_\sigma B})\subset \frac{G\times^BG}{\Ad_\sigma B}$。考虑下图
\begin{small}
\begin{equation}\label{E: cal coh corr to local}
\xymatrix@C=15pt{
\on{pt}&\ar[l] G/B\ar^-{\Delta}[r]& G/B\times G/B\ar@{}[drr]|{\mathrm{(A)}}&& G/B\times G/B\ar_-{\on{sw}(\id\times \sigma)}[ll]\ar@{}[dr]|{\mathrm{(B)}}&\ar_-\Delta[l] G/B\ar[r]& \on{pt}\\
\on{pt}\ar[d]\ar@{=}[u]&\ar[l] \mZ_1\ar[r]\ar[d]\ar[u]& G\times G\ar[d]\ar[u]\ar@{}[drr]|{\mathrm{(C)}}&&\ar_-{\on{sw}(\id\times \sigma)}[ll]G\times G\ar[d]\ar[u]\ar@{}[dr]|{\mathrm{(D)}}&\ar[l]\mZ_2\ar[r]\ar[d]\ar[u]&\on{pt}\ar[d]\ar@{=}[u]\\
\frac{B}{\Ad_\sigma B} & \mZ\ar[l]\ar[r]& \frac{G\times^BG}{\Ad_\sigma B}&&\ar_{\pFr}[ll]\frac{G\times^BG}{\Ad_\sigma B}&\ar[l] \mZ\ar[r]& \frac{B}{\Ad_\sigma B},
}
\end{equation}
\end{small}
其中
\begin{itemize}
\item 方块 $(C)$ 中两个竖直映射是自然投影，映射 $G\times G\to G\times G$ 由 $(g_1, g_2)\mapsto (\sigma(g_2), g_1)$ 给出，使得 $(C)$ 是笛卡儿的；
\item 簇 $\mZ_2$ 的定义使得图 $(D)$ 是笛卡儿的，即 $\mZ_2=\{(g_1,g_2)\in G\times G,\ g_1g_2\in B\}$；
\item 图 $(A)$ 中左边的竖直映射 $G\times G\to G/B\times G/B$ 由 $(g_1,g_2)\mapsto (g_1B, g_2^{-1}B)$ 给出，右边竖直映射由 $(g_1,g_2)\mapsto (g_1^{-1}B,g_2B)$ 给出，使得 $(A)$ 与 $(B)$ 都是笛卡儿的；
\item 簇 $\mZ_1=G$，映射 $\mZ_1\to G\times G$ 由 $g\mapsto (g,g^{-1})$ 给出，使得复合 $\mZ_1\to G\times G\to \frac{G\times^BG}{\Ad_\sigma B}$ 经过 $\mZ$ 分解，且所得态射 $\mZ_1 \to \mZ$ 光滑；而映射 $\mZ_1\to G/B$ 是自然投影（因而也光滑）。
\end{itemize}
把每一行取纤维积，得到交换图
\begin{equation*}
\xymatrix{
\on{pt}&\ar[l] (G/B)(\bF_q)\ar[r] & \on{pt}\\
\on{pt}\ar@{=}[u]\ar[d]&\ar[l] \mW \ar[u]\ar[d]\ar[r] & \on{pt}\ar@{=}[u]\ar[d]\\
B(\bF_q)\backslash \on{pt}&\ar[l] B(\bF_q)\backslash G(\bF_q)/B(\bF_q)\ar[r]& \on{pt}/B(\bF_q),
}\end{equation*}
其中 $\mW:=\mZ_1\times_{G\times G, \on{sw}(\id\times \sigma)}\mZ_2\cong\{g\in G\mid \sigma(g)^{-1}g\in B\}$；它通过 $g\mapsto\sigma(g)$ 映入 $\mZ_1=G$，通过 $g\mapsto(\sigma(g)^{-1},g)$ 映入 $\mZ_2$。映射 $\mW\to (G/B)(\bF_q)$ 把 $g$ 送到 $gB$，故其纤维都同构于 $B$。映射 $\mW\to B(\bF_q)\backslash G(\bF_q)/B(\bF_q)$ 把 $g$ 送到它在 $B(\bF_q)\backslash (G/B)$ 中的像（该像自然落在 $B(\bF_q)\backslash G(\bF_q)/B(\bF_q)$ 中）。

按定义，$S_{\mF}: (B(\bF_q)\backslash \on{pt},\widetilde{\delta_e})\to (\on{pt}/B(\bF_q),\widetilde{\delta_e})$ 是支撑在上图底行的上同调对应。另一方面，有支撑在顶行的上同调对应
$$S'_\mF:=\on{ev}_{\mF}\circ \Ga_{\on{sw}(\id\times\sigma)}^* \circ\on{coev}_{\mF}: (\on{pt},E)\to(\on{pt},E),$$
按 Braverman–Varshavsky 的迹公式（见 \cite[引理 A.2.22]{XZ}），它由 $B(\bF_q)\backslash G(\bF_q)/B(\bF_q)$ 上的函数 $f_\mF$ 给出。

由于 $\mW\to (G/B)(\bF_q)$ 与 $\mW\to B(\bF_q)\backslash G(\bF_q)/B(\bF_q)$ 都是等维光滑的（维数为 $\dim B$），可以用 \eqref{ASS:smooth pullback correspondence}--\eqref{ASS:smooth pullback correspondence3} 中的形式体系把 $S_\mF$ 与 $S'_\mF$ 拉回到中间一行。只须证明它们的拉回相等。

注意 $S_{\mF}$ 与 $S'_{\mF}$ 都是三个上同调对应的复合。
由于交换方块 (A)–(D) 都是笛卡儿的，引理 \ref{AL:pullback compatible with composition} 说明光滑拉回与上同调对应的复合可交换。因此只剩验证在图 \eqref{E: cal coh corr to local} 中
\begin{itemize}
\item 支撑在 $\mZ$ 上的上同调对应 $\widetilde{\coev_\mF}: (\frac{B}{\Ad_\sigma B}, \widetilde{\delta_e})\to (\frac{G\times^BG}{\Ad_\sigma B}, \widetilde{\mF}\tilde\times \widetilde{\mF^*})$ 到中间一行的拉回，等于支撑在 $G/B$ 上的 $\on{un}_{\mF}: (\on{pt},E)\to (G/B\times G/B, \mF\boxtimes \bD\mF)$ 的拉回；关于 $\ev_\mF$ 与 $\on{coun}_{\mF}$ 的对偶论断同样成立；
\item 对应 $\Ga_{\pFr}^*: (\frac{G\times^BG}{\Ad_\sigma B}, \widetilde{\mF}\tilde\times\widetilde{\mF^*})\to (\frac{G\times^BG}{\Ad_\sigma B}, \widetilde{\mF^*}\tilde\times\widetilde{\mF})$ 到中间一行的拉回，等于 $\Ga_{\on{sw}(\id\times\sigma)}^*: (G/B\times G/B, \mF\boxtimes \bD \mF)\to (G/B\times G/B, \bD \mF\boxtimes \mF)$ 的拉回。
\end{itemize}
\end{proof}
